% !TEX program = xelatex
% =============================================================================
% CHP SÖYLEM ANALİZİ RAPORU
%
% Derleme:  xelatex main.tex   (iki kez, içindekiler için)
% ya da:    latexmk -xelatex main.tex
% fontspec kullanıldığı için pdflatex ile derlenmez.
% =============================================================================

\documentclass[12pt,a4paper]{article}

% --- kapak ve sayfa üstü bilgileri ------------------------------------------
\newcommand{\reporttitle}{Hesaplamalı \\ Söylem Analizi Raporu}
\newcommand{\reportsubtitle}{Cumhuriyet Halk Partisi Genel Başkan Konuşmaları, 2015--2026}
\newcommand{\reportshorttitle}{CHP Söylem Analizi}
\newcommand{\reportauthor}{Eren Işık Avcı}
\newcommand{\reportdate}{\today}

\usepackage{preamble}
\usepackage{hyperref}
\usepackage{threeparttable}

\begin{document}

% =============================================================================
% BAŞLIK SAYFASI
% =============================================================================
\thispagestyle{empty}
\vspace*{4cm}

{\color{muted}\large \reportdate}

\vspace{0.7em}
{\fontsize{28}{32}\selectfont\bfseries \reporttitle\par}

\vspace{0.5em}
{\Large\color{muted} \reportsubtitle\par}

\vspace{1.3em}
{\color{ink}\rule{5cm}{2pt}}

\vfill

{\large\bfseries \reportauthor}\\[0.3em]
{\color{muted}
MSc Applied Social Data Science, London School of Economics\\
\href{mailto:erenisikavci@gmail.com}{erenisikavci@gmail.com}
}

\vspace{2cm}

% =============================================================================
% İÇİNDEKİLER
% =============================================================================
\newpage
\thispagestyle{empty}
\tableofcontents

% =============================================================================
% ÖZET


% =============================================================================
% BÖLÜMLER
% =============================================================================
\newpage
\setcounter{page}{1}

\section*{Özet}
\addcontentsline{toc}{section}{Özet}

bulgular ve öneriler

% BLUF: sonuc en basta. Bu sayfa tek basina okundugunda is gormeli.
% Rapor burada bitse bile okuyucu ne yapacagini bilmeli.

\section{Giriş}







% --- kronolojik govde: her olay kendi bolumunde ---

\section{Veri ve Kapsam}

Bu rapor Cumhuriyet Halk Partisi'nin 2015 ve 2026 yılları arasında Kemal 
Kılıçdaroğlu ve Özgür Özel tarafından yapılan konuşmalarını (n = 2.703\footnote{Ham konuşma sayısı = 2.744}) iki farklı nicel metin analizi yaklaşımıyla analiz etmektedir. 
Sözcük düzeyindeki analizde 6 farklı sözlüğün 9 farklı kırılma anında nasıl davrandığı,
bağlam düzeyindeki analizdeyse 11 yıl boyunca \textit{millet}, \textit{kadın} ve \textit{erdoğan} kelimelerinin
anlam kümeleri incelenmiştir.

Raporda analiz edilen konuşmalar \href{https://chp.org.tr}{CHP'nin internet sitesinden} otomatik olarak çekilmiş ve
veri toplama esnasında araştırmanın amacı\footnote{London School of Economics and Political Science (LSE) Metodoloji Departmanı, Uygulamalı Sosyal Veri Bilimi yüksek lisans programı bitirme tezi.} user-agent kullanılarak tanımlanmıştır. Ayrıca her toplanan metin arasında
sunucuya yük bindirmemek için makul (5 saniye) bir gecikme süresi uygulanmıştır. 

Korpus, etkinlik (873), basın (800), miting (445), meclis (347), televizyon (214), ve kurultay (24)
olmak üzere 6 farklı mekandaki konuşmalardan bir araya getirilmiştir\footnote{Birden fazla başlıkta paylaşılan konuşmalar ayıklandı}.
Miting konuşmalarının coğrafi dağılımı [TABLO X]'te paylaşılmıştır.
13 Ocak 2015 ve 8 Kasım 2023 arasında Kemal Kılıçdaroğlu'nun 1.569; 9 Kasım 2023 ve 31 Mart 2026 arasında Özgür Özel'in 1.134 konuşması tespit edilmiştir. 
Görev süresince Özgür Özel daha fazla basın ve miting; Kemal Kılıçdaroğlu ise daha fazla meclis ve etkinlik konuşması gerçekleştirmiştir.
Mekan yoğunluğunun hem sözcük hem bağlamsal bulguları etkileyip etkilemediğine yönelik dayanıklılık testleri yapılmış ve mekanın bulguları etkilemediği görülmüştür.
Bulguların tespit edilemeyen bir kısmı Özgür Özel ve Kemal Kılıçdaroğlu arasındaki Türkçe kullanımından kaynaklanıyor olabilir. Örneğin, Kılıçdaroğlu'na kıyasla Özel daha
çeşitli bir kelime dağarcığı tercih etmekte ve cümle başına yaklaşık 2 kelime daha fazla kullanmaktadır. Buna ek olarak rapor boyunca liderlerin retorik kullanımlarına yönelik bulgulara da yer verilmiştir.


Bu araştırmanın kapsamında yalnızca parti liderlerinin konuşmalarına odaklanılmıştır. Milletvekillerinin açıklamaları için
partinin internet sitesinde seçilmiş konuşmalar yayınlanmakta olup tam kapsamlı bir korpus için Türkiye Büyük Millet Meclisi arşivine başvurulabilir.
Ekrem İmamoğlu'nun konuşmalarının analiz edilmesi planlansa da yeterli veri kaynağına erişilememiştir. Liderlerin sosyal medya 
paylaşımları internet sitesinden erişilebilir olsa da yeterli hacimde bir veri kaynağı oluşturmadığından analiz dışında bırakılmıştır.

Dokümanların nasıl temsil edildiği, korpus temizlik aşamaları, kullanılan yöntemler ve algoritmalar yöntem bölümünde açıklanmıştır. 
Frekans ve kırılma analizi için oluşturulan sözlükler yöntem bölümünde mevcuttur.
Analizi yeniden üretmek veya kodu ve dayanıklılık testlerini incelemek isteyen araştırmacılar e.avci@lse.ac.uk adresinden iletişime geçebilir.





% =============================================================================
% Ekonomik Talepler
% =============================================================================

\section{Ekonomik Talepler}
Ekonomik talepler kırılma anlarına duyarlı. Geri dönüş 4 ila 7 ay sürüyor.
En yüksek artış kur krizini de kapsayan 2018 Genel Seçimlerinin ardındaki pencerede gözlemlendi. 

\begin{table}[H]
\centering
\caption{Trendlerine göre ekonomik kayıtlar}
\label{tab:olay-kayit}
\small
\begin{tabular}{@{}l l r r r c l@{}}
\toprule
Olay & Taban & Sonrası & Değişim & & Dönüş \\
\midrule

% ---------------------------------------------------------------- ÇEKİLENLER
\multicolumn{6}{@{}l}{\textbf{Çekilenler}} \\[0.2em]
  Darbe girişimi & 74.5 & 42.3 & $-$43.2\% & \azalan & 4 ay \\
   Anayasa referandumu & 55.7 & 34.8 & $-$37.5\% & \azalan & 4 ay \\
   Yerel seçim 2019    & 174.9 & 112.9 & $-$35.4\% & \azalan & 6 ay \\
   İmamoğlu tutuklanması & 143.4 & 95.1 & $-$33.7\% & \azalan & 7 ay\\

\addlinespace[0.6em]
% ---------------------------------------------------------------- YÜKSELENLER
\multicolumn{6}{@{}l}{\textbf{Yükselenler}} \\[0.2em]
   Genel seçim 2018    & 109.4 & 200.4 & $+$83.2\% & \artan & \\ % dipnot: yükselmenin temel sebebi kur krizi
      Bahçeli'nin çağrısı & 129.0 & 190.5 & $+$47.7\% & \artan & \\
   Genel seçim 2023 & 115.3 & 144.8 & $+$25.6\% & \artan & \\
      CHP Istanbul kayyum & 125.2 & 140.4 & $+$12.2\% & \artan & \\

\addlinespace[0.6em]
% ------------------------------------------------------------ HAREKET ETMEYEN
\multicolumn{6}{@{}l}{\textbf{Hareket etmeyenler}} \\[0.2em]
Yerel seçim 2024              & 124.6 & 129.6 & $\pm$4.0\% & \duraga & \\

\bottomrule
\end{tabular}
\kaynaknot{Taban olaydan önceki 12 ayın, \emph{Sonrası} olaydan sonraki
üç ayın ortalamasıdır; ikisi de 10.000 kelimede vukuat sayısıdır, yani
konuşma hacminden bağımsızdır. \emph{Dönüş}, kaydın tabanın \%90'ına geri
döndüğü ay sayısıdır. Korpus 31 Mart 2026'da bittiği için Bahçeli çağrısı,
İmamoğlu tutuklaması ve İstanbul kayyumu sonrasında 24 aylık pencere
doldurulamıyor; oralarda dönüş \emph{gözlem dışıdır}, gerçekleşmediği
anlamına gelmez.}
\end{table}



% =============================================================================
% REFERANDUM, Nisan 2017
% =============================================================================
\newpage
\section{Kimlikler}

Hem ideolojik hem de etnik ve inanç kimlikleri olaylara duyarlı. 
Etnik ve inanç kimlikleri Özgür Özel döneminde söylemde yer ediniyor. CHP söylemi ideolojik kimliklerden çok etnik
ve inanç kimliklerini yöneliyor. Etnik ve inanç kimlikleri bütün olaylara tepki gösterirken ideolojik kimlikler referandum, 2019 yerel seçimleri ve İmamoğlu'nun
tutuklanması gibi dislokasyonlara duyarsız kalıyor. 


\begin{table}[H]
\centering
\caption{Trendlerine göre ideolojik ile etnik ve inanç kimlikleri kayıtları}
\label{tab:olay-kayit}
\small
\begin{tabular}{@{}l l r r r c l@{}}
\toprule
Kayıt & Olay & Taban & Sonrası & Değişim & & Dönüş \\
\midrule

% ---------------------------------------------------------------- ÇEKİLENLER
\multicolumn{7}{@{}l}{\textbf{Çekilenler}} \\[0.2em]
Etnik ve inanç & Darbe girişimi, 2016 & 2.3 & 0.6 & $-$72\% & \azalan & > 24 ay \\
Etnik ve inanç & Referandum, 2017 & 0.6 & 0.4 & $-$42\% & \azalan & 4 ay \\
Etnik ve inanç & Genel seçim, 2018 & 1.0 & 0.0 & $-$95\% & \azalan & 15 ay \\
İdeolojik & Genel seçim, 2018 & 4.9 & 3.3 & $-$33\% & \azalan & 3 ay \\
Etnik ve inanç & Genel seçim, 2023 & 1.1 & 0.2 & $-$85\% & \azalan & 5 ay \\
İdeolojik & Bahçeli çağrısı, 2024 & 7.0 & 4.7 & $-$32\% & \azalan & 4 ay \\
İdeolojik & İstanbul kayyumu, 2025 & 5.2 & 2.8 & $-$46\% & \azalan & gözlem dışı \\

\addlinespace[0.6em]
% ---------------------------------------------------------------- YÜKSELENLER
\multicolumn{7}{@{}l}{\textbf{Yükselenler}} \\[0.2em]
İdeolojik & Darbe girişimi, 2016 & 0.6 & 2.3 & $+$283\% & \artan \\
Etnik ve inanç & Yerel seçim, 2019 & 0.3 & 0.8 & $+$142\% & \artan \\
Etnik ve inanç & Yerel seçim, 2024 & 2.3 & 5.6 & $+$138\% & \artan \\
Etnik ve inanç & Bahçeli çağrısı, 2024 & 5.1 & 6.4 & $+25$\% & \artan \\
Etnik ve inanç & İmamoğlu tutuklanması, 2025 & 6.0 & 9.8 & $+$62\% & \artan \\
Etnik ve inanç & İstanbul kayyumu, 2025 & 8.0 & 8.9 & $+$12\% & \artan \\


\addlinespace[0.6em]
% ------------------------------------------------------------ HAREKET ETMEYEN
\multicolumn{7}{@{}l}{\textbf{Hareket etmeyenler}} \\[0.2em]
İdeolojik & Referandum, 2017 & 3.6 & 4.0 & $+$10\% &  &  \\
İdeolojik & Yerel seçim, 2019 & 5.1 & 5.1 & 0\% & & \\
İdeolojik & İmamoğlu tutuklaması, 2025 & 5.8 & 6.0 & $+$3\% & & \\

\bottomrule
\end{tabular}
\kaynaknot{Taban olaydan önceki 12 ayın, \emph{Sonrası} olaydan sonraki
üç ayın ortalamasıdır; ikisi de 10.000 kelimede vukuat sayısıdır, yani
konuşma hacminden bağımsızdır. \emph{Dönüş}, kaydın tabanın \%90'ına geri
döndüğü ay sayısıdır. Korpus 31 Mart 2026'da bittiği için Bahçeli çağrısı,
İmamoğlu tutuklaması ve İstanbul kayyumu sonrasında 24 aylık pencere
doldurulamıyor; oralarda dönüş \emph{gözlem dışıdır}, gerçekleşmediği
anlamına gelmez.}
\end{table}

\subsection{Kadın: Semantik}
\begin{figure}[H]
    \centering
    \includegraphics[width=\textwidth]{kadın_phase_heatmap.pdf}
    \caption[Phase heatmap: \emph{kadın}]{Kadın kelimesi etrafındaki bağlam kümelerinin
    faz oranı. Oran her fazdaki bir kümenin payının tüm korpustaki payına oranıdır. 
    1.0 değeri orantılı varlığı; 1.0 üzeri veya altı değerler göreli yoğunluk veya kıtlığı gösterir.
    Renk koyulaştıkça ilgili kümenin o fazdaki yoğunluğu artar. $KK_{Erken} = 13.01.2015-01.04.2019$; 
    $KK_{Gec} = 02.04.2019-09.11.2023$; $Ozel = 06.11.2023-31.03.2026$. Sırasıyla dönem frekansları: 1: 1161, 2: 1884, 3: 2563.}
    \label{fig:heatmap-kadın}
\end{figure}

Kılıçdaroğlu, erken dönemde kadınlara yönelik söylemini popüler taleplerle birlikte dile getirdiği kadın-erkek eşitliği
ve genç, yaşlı, erkek ve çocuğun beraber-kapsayıcı mücadelesi üzerinden inşa etmektedir. Okuryazarlık ve Eğitim ile Seçme
ve Seçilme kümeleri bir arada değerlendirildiğinde Kemal Kılıçdaroğlu erken Cumhuriyet döneminin 
kazanımlarına daha yoğun atıfta bulunmaktadır. Bahşedilen Haklar kümesinde Kılıçdaroğlu,
"Cennet annelerin ayakları altındadır." hadisi ve çeşitli dini referanslar\footnote{Benzeri retorik dini referanslara tezimde "hak" kelimesinin etrafındaki semantik kümeleri incelerken de rastladım. Örn: Kul hakkı, yetimin ve öksüzün hakkı gibi tamlamalar.}
üzerinden kadına şiddeti ve AKP'nın önlemsizliğini eleştirmekte, aynı kümede
Özgür Özel ise dini referansları bırakıp yerine, Kılıçdaroğlu'nda da gözlemlendiği gibi, erken Cumhuriyet döneminde kadınlara tanınan haklardan söz etmektedir.

Geç dönem söz konusu olduğunda Kılıçdaroğlu'nun söylemindeki Atatürk döneminin getirilerine yönelik anlatı
sönümlenmiş ve önceki dönemden aktif olan 8 kümeden yalnızca Cinsiyet Eşitliği, kapsayıcı söylem, Kadınların Siyasete Katılımı ve ahlaki 
retoriği devam etmiştir. Bu dönemde Kılıçdaroğlu'nun sosyal yardımlara yönelik vaatleri ilgi çekici olup (örn. Aile Destekleri Sigortası ve Sosyal Güvenlik Primi\footnote{Bu vaatten ilgili kümede ilk kez 3 Eylül 2021 tarihinde bahsediyor.}) ağırlıklı olarak 2023 Mayıs ayındaki mitinglerde yoğunlaşmaktadır.
2020 yılının Temmuz ayından başlamak itibariyle CHP lidelerinin söylemine kadın örgütleriyle yapılan
görüşmeler de konu olmaya başlamıştır. Kılıçdaroğlu, 306 kadın örgütüyle yazılmış yeni siyasi partiler yasasını ve bu yasadaki
50-50 kadın kota önerisini sıklıkla dile getirirken Özgür Özel kadına yönelik şiddet, 6284 yasası ve Etki Ajanlığı Yasası hakkında
eleştiriler gündeme getirmektedir.

Isı haritasının orta kısmında yer alan kümeler CHP'nin son 11 yılda kadınlar hakkındaki
söyleminin temel bileşenleridir. Bu kümeler sırasıyla şu şekildedir: Kadın tutuklular ve mağdurlar, İstanbul Sözleşmesi ve kadına yönelik şiddet,
CHP belediyelerinin kreş girişimleri ve sosyal adaletsizlik, ve sosyal kimlikler. "Mesleki Eşdeğerlilik" anlatısı altında Kılıçdaroğlu işçi, çiftçi, emekli ve esnaf gibi 
"çalışan" sınıfları hem ekonomik problemlerle yüzleşen hem de ortak bir dayanışma içinde olan "sosyal kimlikler" olarak eklemlemektedir.
Ayrıca bütün bu kimliklerin çözümünün CHP'de olduğu, Bay Kemal'in bu kimlikler için çalışacağı gibi retorik izler de kümede mevcuttur.

Her ne kadar ağırlığı Kılıçdaroğlu'nun ikinci döneminde olsa da 
"Çalışma" kümesi de bu listeye dahil edilebilir. Kılıçdaroğlu, özellikle kırsaldaki kadınların ev içi emeğinin karşılıksız kalmasını eleştiriyorken
Özgür Özel, İBB tarafından açılan kreşlerin kadınların sosyal hayatına etkisine değinmektedir.

Özel döneminde partinin kadın söylemi AKP ve Erdoğan'a eklemlenmeye başlamıştır. Sırasıyla, kabinedeki tek bakanın Aile ve Sosyal Hizmetler
Bakanı olması, Erdoğan ailesine yönelik İran rejimindeki "kadın haklarına" yönelik ihlaller, Filistin'de kadın ve çocuklara yönelik uygulanan katliam,
Suriyeli kadın göçmenlerin yaşadığı zorluklar, MHP ve AKP'nin HÜDAPAR ile ittifak yapması kadınlara yönelik son yıllardaki CHP gündeminin
temel gündemlerini oluşturmaktadır. Buna ek olarak Özgür Özel döneminde gençler ve kadınlara yönelik ortak siyasa önerileri (örn: kota) 
söylemde daha sık yer edinmeye başlamıştır. İlginç bir biçimde Özgür Özel AKP'ye oy veren kadın seçmenlerin çocuklarının aldığı eğitimden memnuniyet oranı Özgür Özel'in
söyleminde sürekli gündeme gelmektedir. 2024 Eylül ayından başlayarak 2026 Ocak ayının başına kadar bu verinin \%17 ila \%21 arasında olduğu tarafınca yinelenmiştir\footnote{Bu veriyi nereden aldınız, güvenemedim.}. 

Özetle, 
\vspace{-1.2em}
\begin{itemize}
    \item Özgür Özel döneminde CHP'nin kadın söylemi daha politize olmuştur. Bunun temel sebebi Erdoğan'a, ailesine, AKP'ye ve bakanlığa yönelik direkt ve 
    kişisel eleştiriler, yani karşı bir düşman inşası, söyleme girmesidir\footnote{Yalnızca "kadın" etrafındaki semantik analizden çıkarılan sonuçtur, diğer kümeler ve frekanslar beraber değerlendirilmeden böyle bir sonuca varılamaz.}. Kemal Kılıçdaroğlu'nun kadınlara yönelik söylemi sosyal yardımlara öncelik veren ideal bir anlatıya bağlı kalırken
    Özgür Özel işsizlik, yoksuzlluk, kreş, kota uygulamaları gibi güncel ve siyasa üretmeye yönelik talepleri vurgulamaktadır.
    \item 
\end{itemize}

\subsection{Genç: Semantik}

% burada "zincir" kullandım, kadın kümelerindeki çalışma için de Zincir: Çalışma yazabilirim
 

% =============================================================================
% 2023-2026
% =============================================================================

\section{Antagonist}

% Bu bolumde birden fazla olay var: 2023 genel secimi, lider degisimi,
% 2024 yerel secimi, Bahceli cagrisi, Imamoglu tutuklamasi, Istanbul kayyumu.
%
% Kritik ayrim, ikisinin AYNI ANDA olmadigi:
%   Yolsuzluk ve cete   kirilma 2023-11  (lider degisimiyle ayni ay)
%   Erdogan ve AKP      kirilma 2024-09  (on ay sonra)
% Saf uslup etkisi olsaydi ikisi de 2023-11'de kirilirdi.
%
% Son uc olayda korpus 24 ayi doldurmuyor: donus GOZLEM DISI, "gerceklesmedi"
% degil.

\subsection*{Erdoğan ve AKP}

2023 Mayıs genel seçimlerinden sonraki dönemde Erdoğan ve AKP kelimelerinin frekansları
sürekli olarak artışa geçiyor ve İmamoğlu tutuklanmasından sonraki üç aylık periyotta pik yapıyor.
Buna karşı olarak darbe girişimi ve referandum sonrasında aynı tema sönümleniyor.
Düşüş oranları benzer olsa da eski seviyeye dönüş darbe girişiminde 16, referandumda 5 ay sürüyor.

\begin{table}[H]
\centering
\caption{Erdoğan ve AKP: olay çevresindeki hareket}
\label{tab:olay-antagonist_akponly}
\small
\begin{tabular}{@{}l r r r c l@{}}
\toprule
Olay & Taban & Sonrası & Değişim & & Dönüş \\
\midrule
  Darbe girişimi, Tem. 2016 & 16.5 & 9.7 & -41\% & \azalan & 16 ay \\
  Referandum, Nis. 2017 & 11.2 & 7.1 & -37\% & \azalan & 5 ay \\
  Genel seçim, Haz. 2018 & 12.6 & 14.9 & +18\% & \artan & --\\
  Yerel seçim, Mar. 2019 & 15.2 & 24.1 & +58\% & \artan & -- \\
  Genel seçim, May. 2023 & 9.9 & 18.8 & +90\% & \artan & -- \\
  Yerel seçim, Mar. 2024 & 21.9 & 31.8 & +45\% & \artan & -- \\
  Bahçeli çağrısı, Eki. 2024 & 31.4 & 43.0 & +37\% & \artan & -- \\
  İmamoğlu tutuklaması, Mar. 2025 & 36.3 & 53.1 & +46\% & \artan & -- \\
  İstanbul kayyumu, Eyl. 2025 & 49.7 & 50.0 & +1\% & \duraga & -- \\
\bottomrule
\end{tabular}
\kaynaknot{Taban olaydan önceki 12 ayın, \emph{Sonrası} sonraki üç ayın
ortalamasıdır; ikisi de 10.000 kelimede vukuat sayısıdır. \emph{Dönüş},
kaydın tabanın \%90'ına döndüğü ay sayısıdır. Korpus 31 Mart 2026'da bittiği
için son üç olayda 24 aylık pencere dolmuyor: oralarda dönüş gözlem dışıdır,
gerçekleşmediği anlamına gelmez.}
\end{table}


\subsection*{Yolsuzluk ve Çete}

Kılıçdaroğlu döneminde anayasa referandumundan sonra tırmanmaya başlayan yolsuzluk ve çete teması 
İBB seçimlerinin yenilendiği 2019 yerel seçimleri sonrasındaki 3 ay pencere içerisinde $-$46\% oranında azalıyor.
13 ay sonra eski seviyesine dönse de çete, mafya, baron ve rüşvet gibi kelimeler lider değişimiyle birlikte 
söylemden büyük oranda çekiliyor ($-$66\%). [ÖZEL DÖNEMİNDEKİ İKİ ARTIŞ MUHTEMELEN ADALET BAKANLIĞINDAN
KAYNAKLANIYOR. AMA GÖRMEDEN GEÇEMEM. VERİ İLE KONTROL EDİP BURAYA DÖN]




\begin{table}[H]
\centering
\caption{Yolsuzluk ve çete: olay çevresindeki hareket}
\label{tab:olay-antagonist_general}
\small
\begin{tabular}{@{}l r r r c l@{}}
\toprule
Olay & Taban & Sonrası & Değişim & & Dönüş \\
\midrule
  Darbe girişimi, Tem. 2016 & 3.9 & 2.5 & -34\% & \azalan & 10 ay \\
  Referandum, Nis. 2017 & 2.0 & 4.5 & +128\% & \artan & -- \\
  Genel seçim, Haz. 2018 & 11.8 & 29.6 & +151\% & \artan & -- \\
  Yerel seçim, Mar. 2019 & 21.4 & 11.6 & -46\% & \azalan & 13 ay \\
  Genel seçim, May. 2023 & 25.0 & 22.5 & -10\% & \duraga & -- \\
  Yerel seçim, Mar. 2024 & 14.9 & 4.4 & -71\% & \azalan & >24 ay \\
  Bahçeli çağrısı, Eki. 2024 & 4.7 & 3.4 & -27\% & \azalan & 4 ay \\
  İmamoğlu tutuklaması, Mar. 2025 & 3.8 & 5.5 & +47\% & \artan & -- \\
  İstanbul kayyumu, Eyl. 2025 & 5.7 & 7.8 & +36\% & \artan & -- \\
\bottomrule
\end{tabular}
\kaynaknot{Taban olaydan önceki 12 ayın, \emph{Sonrası} sonraki üç ayın
ortalamasıdır; ikisi de 10.000 kelimede vukuat sayısıdır. \emph{Dönüş},
kaydın tabanın \%90'ına döndüğü ay sayısıdır. Korpus 31 Mart 2026'da bittiği
için son üç olayda 24 aylık pencere dolmuyor: oralarda dönüş gözlem dışıdır,
gerçekleşmediği anlamına gelmez.}
\end{table}




% --- olaylar arasi duzenlilikler ---
% =============================================================================
% SENTEZ
% =============================================================================

\section{Milleti İnşa Etmek}

\newpage

% =============================================================================
% SENTEZ MALZEMESI  (analizden cikan, tek bolumde soylenemeyen bulgular)
%
% Ayrimin kurali: bu cumleyi TEK BIR DILE bakarak soyleyebiliyor muyum?
% Evetse o dilin bolumune ait. Hayirsa buraya.
%
% -----------------------------------------------------------------------------
% 1. IKAME. Iki kayit ters yone gidiyor ama AYNI ANDA DEGIL.
%      Yolsuzluk ve cete   kirilma 2023-11   -66%  (lider degisimiyle ayni ay)
%      Erdogan ve AKP      kirilma 2024-09  +169%  (on ay sonra)
%    Ikisi de q<0.001, aylik seride, basamak (kayma degil).
%    Yani "yeni lider yeni dil" degil: eski dil liderle birakildi, yeni dil
%    on ay sonra ayri bir konjonkturde basladi. Iki ayri karar.
%
% 2. KONFAUNDA KARSI KANIT. Hareket etmeyenler hareket edenler kadar onemli:
%      Ekonomik talepler   +3%
%      Kadin (anma)        -0%
%    Genel bir uslup kaymasi olsaydi hepsi birlikte oynardi. Yolsuzluk %64
%    duserken ekonomik dilin hic kimildamamasi, degisimin SECICI oldugunu
%    gosteriyor. Bu ancak oynayanlarla oynamayanlar yan yana konunca gorunur.
%
% 3. MEKAN BULGULARI TASIMIYOR. Son tabakalama ile mekan bilesimi iki donemde
%    esitlendi (agirliklar: meclis %25, etkinlik %24, tv %20, miting %16,
%    basin %14, kurultay %1). On uc kayittan on birinde sonuc degismedi:
%      Erdogan ve AKP     +176% -> +177%
%      Yolsuzluk ve cete   -64% ->  -64%
%      Ahlaki hak edis     -77% ->  -78%
%    Iki istisna (Ideolojik kimlikler, Kadin haklari) ve ikisinde de duzeltme
%    etkiyi BUYUTUYOR. Yani mekan bulgu yaratmiyor, gizliyor.
%    Bu tek bolumde soylenirse on uc kez tekrarlanir, burada bir kez.
%
%
% -----------------------------------------------------------------------------
% YAPISAL NOT
% Bu bolumun adi su an "Bulgular" ama ondan onceki bes bolum de bulgu.
% Ad, bolumun ne YAPTIGINI soylemeli: "Diller Arasinda Tekrar Edenler" gibi.
% Yeri de leksikal blok ile semantik blok arasinda kaliyor; bir sentez
% sentezledigi seyden sonra gelmeli, yani 09_erdogan'in ardina.
%
% Asagidaki birinci tablo tamamen "Ekonomik talepler" satirlarindan olusuyor,
% yani tek dilin olay davranisi: 02_ekonomik'e ait, buraya degil.
% =============================================================================

\section{Bulgular}

% Her bulgunun basligi IDDIANIN KENDISI. Basliklar icindekiler listesine
% dusuyor, yani okuyucu sadece icindekilere bakarak argumani gorebiliyor.

\begin{bulgu}{Yolsuzluğa ve çeteye karşı ahlaki çerçeve sönümleniyor. Etnik ve inanç kimlikleri
söylemde yer ediniyor.}

Kul hakkı gasp edilmiş ve hak, hukuk, adalet mücadelesi veren insanlara karşı beşli çete, mafya ve baron
gibi figürlerle doldurulan düşman anlatısı Özgür Özel döneminde terk ediliyor. 
Erdoğan ve AKP kelimelerinin frekansı 2024 Eylül'de 2.7 kat artıyor. Etnik ve inanç kimlikleri
2024 yerel seçimlerinden sonra CHP söyleminde düzenli bir biçimde çağrılmaya başlanıyor. 


\begin{table}[H]
\centering
\caption{Kayıtların iki dönem arasındaki değişimi}
\label{tab:register-change}
\begin{tabular}{@{}l r r r c@{}}
\toprule
 & Kılıçdaroğlu & Özel & Değişim (kat) & \\
\midrule
Etnik ve inanç kimlikleri & 0.6 & 7.6 & +12.5x & \artan \\
Erdoğan ve AKP & 16.5 & 45.5 & +2.8x & \artan \\
Ahlaki çerçeve & 11.1 & 2.5 & -4.3x & \azalan \\
Yolsuzluk ve çete & 14.9 & 5.3 & -2.8x & \azalan \\
\bottomrule
\end{tabular}
\kaynaknot{...}
\end{table}

\end{bulgu}

\begin{bulgu}{Ekonomik talepler ve sınıflar söylemin en kırılgan(?) olabileceği anlara duyarlı. İki kategori CHP
  söyleminde sürekli kesintiye uğruyor.} 
Siyasi bir rejimin sıklıkla deneyimlemediği ve tarafların yeni bir söylem geliştirmek zorunda
kaldığı anlar (dislokasyon), diğer temalara kıyasla sistematik bir şekilde ekonomik talepleri ve sınıfları etkiliyor.
İktidarın etkisinin olduğu dislokasyonlarda ekonomik taleplerin eski seviyelerine dönmesi darbe ve anayasa
referandumuna kıyasla 2 ila 3 ay daha uzun sürüyor. CHP söylemi darbe girişiminden sonraki üç aylık süreçte ekonomik talepleri ve rakibe yönelik dili (örn. diktatör) bırakıyor, buna karşılık kuruluş anlatısını ve ideolojik kimlikleri öne çıkarıyor. 
  % burada en kırılgandan kastım dislokasyon olduğunu söylemdeden anlatmak. yani bu anlar
  % söylemin yeniden kurgulanmak zorunda kalındığı anlar



\begin{table}[htbp]
\centering
\small
\setlength{\tabcolsep}{4pt}
\caption{Log-likelihood keyness of CHP discourse around seven critical events (signifier-independent): words that \textsc{enter} (keyness $+$) vs.\ \textsc{recede} (keyness $-$) across the event, in symmetric $\pm 3$-month windows ($\geq 5$ speeches). Events marked $\dagger$ are treated as dislocations. Each event reactivates a \emph{distinct} register rather than following a single law. Lemma artefacts merged (\textit{üre}$\rightarrow$\textit{üretim}); place-names omitted.}
\label{tab:dislocation-all}
\begin{threeparttable}
\begin{tabular}{@{}>{\raggedright\arraybackslash}p{2.6cm} >{\raggedright\arraybackslash}p{5.3cm} >{\raggedright\arraybackslash}p{5.3cm}@{}}
\toprule
\textbf{Dislokasyon} & \textbf{Enters (reactivated)} & \textbf{Recedes} \\
\midrule
Darbe girişimi†, 2016 & darbe\tnote{a}, OHAL, liyakat, cumhuriyet, laik & dokunulmazlık, üretim\tnote{b}, asgari\_ücret, lira, elektrik, \textbf{diktatör}, turizm, [sosyal] yardım \\
\addlinespace
Referandum†, 2017 & adalet, darbe, savcı, yürüyüş, gayrımeşru, meşru, saray & anayasa, hayır, millet, ulusal, \textbf{tek\_adam\_rejimi}, otorite, Suriye \\
\addlinespace
Genel seçimler, 2018 & \textbf{kriz}, ihale, merkez\_bankası, israf, dolar, borç & esnaf, şeker [fabrika], barışş, kadın, eğitim, emekli, okul, kadın, \textbf{vesayet}, aday \\
\addlinespace
Yerel seçimler, 2019† & adalet, mazbata, hakim, vicdan, iptal & milliyetçi, hizmet, \textbf{özelleştirme}, lira, \textbf{fabrika}, kuyruk, belediye, mahalle, vergi \\
\addlinespace
Genel seçimler, 2023 & \textbf{Erdoğan}, borç, gazeteci, ahlak, dolandırıcı, montaj, kurultay, medya, irade & kul\_hakkı, emekli, genç, beşli\_çete, deprem, kuruş, öğretmen, adalet, okul, vaat kelimeleri\tnote{c} \\
\addlinespace
Yerel seçimler, 2024 & anayasa, birinci, \textbf{Gezi}, \textbf{normalleşme}, müzakere, IMF, müfredat, buğday, kamu & emekli, ittifak, aday, lira, altın, saray, DEM Parti, Erdoğan \\
\addlinespace
İmamoğlu tutuklanması 2025† & İmamoğlu, Saraçhane, sivil\_darbe, cunta, kayyum, boykot, genç, terörsüz, Silivri, rüşvet, diploma & asgari\_ücret, geçim, emekli, maaş, kadın, orman\_yangını, deprem\tnote{d} \\
\bottomrule
\end{tabular}
\begin{tablenotes}\footnotesize
\item[a] Ve FETÖ, cemaat, temmuz ve girişim gibi kelimeler.
\item[b] Ve çiftçi, tarım, fındık ve ürün gibi üretimle ilgili kelimeler.
\item[c] Söz, getireceğim, getireceğiz, yapacağız, göreceksiniz\dots
\item[d] Ve 6 Şubat 2023 depremleriyle ilgili kelimeler: Adıyaman, Hatay, Kahramanmaraş, container housing.
\end{tablenotes}
\end{threeparttable}
\end{table}


\end{bulgu}

\begin{bulgu}{Geçmişin antagonistleri (buraya başlık lazım bir de ne tartışacağım)}

\begin{table}[H]
\centering
\caption{Kayıtların iki dönem arasındaki değişimi}
\label{tab:register-change-akp}
\begin{tabular}{@{}l r r r c@{}}
\toprule
 & Kılıçdaroğlu & Özel & Değişim (kat) & \\
\midrule
Darbe çerçevesi & 10.7 & 16.3 & +1.5x & \artan \\
Cemaat & 5.0 & 1.8 & -2.7x & \azalan \\
Terör & 11.1 & 5.3 & -2.1x & \azalan \\
\bottomrule
\end{tabular}
\kaynaknot{...}
\end{table}

\end{bulgu}

\begin{bulgu}{Emek}

%Emek ve çalışma & 23.1 & 40.9 & +1.8x & \artan \\

\end{bulgu}

% burayı hegemonik analiz olarak yorumlayamam gibime geliyor tartışacağım
%\multicolumn{5}{@{}l}{\textbf{Değişmeyenler}} \\
 % Ekonomik talepler & 129.3 & 133.4 & 1.0x & \duraga \\
  %Kadın hakları & 13.1 & 13.0 & 1.0x & \duraga \\
  %İdeolojik kimlikler & 6.3 & 5.6 & 1.1x & \duraga \\
  %Kurtuluş anlatısı & 5.5 & 4.6 & 1.2x & \duraga \\








\section{Hangi adalet?}

\section{Erdoğan (buraya farklı başlık)}

\subsection{Erdoğan Semantik}
\section{Öneriler}

% Her oneri bir bulguya bagli olmali. Bulgusuz oneri fikirdir, analiz degil.

\begin{oneri}{Sosyal kimlikler hangi tema olursa olsun taleplerden kopmamalı.}
\textbf{Dayanak:} Bulgu 2

[ne yapılacağı, somut]
\end{oneri}

\begin{oneri}{[öneri başlığı]}
\textbf{Dayanak:} Bulgu 2.

[ne yapılacağı]
\end{oneri}

\section{Yöntem}

% En sonda ve kisa. Isteyen buradan okur, cogu okuyucu okumaz.
% Teknik terimler burada kalabilir, ana metinde kalamaz.

\subsection*{Veri}
[kaynak, toplama, temizleme, kapsam]

\subsection*{Sözcük düzeyi}
[olay pencereleri, sözlükler, ağırlıklandırma]

\subsection*{Bağlam düzeyi}
[gömme, kümeleme, kırılma tespiti]

\subsection*{Bu analiz neyi söyleyemez}
[sınırlar. Baştan söylemek, sonradan sorulmasından iyidir.]


\end{document}
